% Source: https://artofproblemsolving.com/community/c3h2173734_finding_the_coefficient_of_a_term_in_a_large_expansion
% Stable identity: aops-topic-2173734-q01
Factor out $x^{28}$:
\[\left(x+x^2+\cdots+x^7\right)^{28}=x^{28}\left(1+x+\cdots+x^6\right)^{28}=x^{28}\left(\dfrac{1-x^7}{1-x}\right)^{28},\]
so we need the coefficient of $x^{69}$ in $\left(1-x^7\right)^{28}(1-x)^{-28}$. (Equivalently: the number of ways to write $69$ as an ordered sum $e_1+\cdots+e_{28}$ with $0\le e_i\le6$.) By the binomial theorem,
\[\left(1-x^7\right)^{28}=\sum_{j=0}^{28}(-1)^j\dbinom{28}{j}x^{7j},\qquad(1-x)^{-28}=\sum_{m\ge0}\dbinom{m+27}{27}x^m .\]
Collecting the terms with $7j+m=69$ (so $0\le j\le9$), the coefficient is
\[\sum_{j=0}^{9}(-1)^j\dbinom{28}{j}\dbinom{96-7j}{27}.\]
This is the inclusion-exclusion count: $\dbinom{96}{27}$ counts all solutions of $e_1+\cdots+e_{28}=69$ in nonnegative integers (stars and bars), and the $j$-th term corrects for $j$ of the $e_i$ being at least $7$ (if $e_i\ge7$, replace $e_i$ by $e_i-7$). So the coefficient is
\[\boxed{\sum_{j=0}^{9}(-1)^j\dbinom{28}{j}\dbinom{96-7j}{27}}.\]
(Evaluating the ten terms, which is a job for a computer, gives $6410611349709577216080$.)
